\documentclass[10pt,a4paper]{amsart}
\usepackage[margin=19mm]{geometry}
\usepackage{amsmath,amssymb,mathtools}
\usepackage[scr=rsfs,cal=euler]{mathalfa}
\usepackage{xfrac}
\usepackage[unicode,hidelinks,bookmarks=false]{hyperref}
\newcommand{\ie}{i.e.\ }
\newcommand{\cf}{cf.\ }
\newcommand{\resp}{respectively\ }
\newcommand{\Subsection}{\subsection}
\providecommand{\nobreakdash}{\nobreak\hbox{-}}
\setlength{\emergencystretch}{2em}
\multlinegap0pt
\usepackage{bm}
\usepackage{bbm}
\usepackage{dsfont}
\usepackage{xspace}
\usepackage{cancel}
\usepackage{mathtools}
\usepackage{cleveref}
\usepackage[shortlabels]{enumitem}
\usepackage{amsthm}
\makeatletter
\@ifundefined{teo}{\newtheorem{teo}{Teo}}{}
\@ifundefined{lemma}{\newtheorem{lemma}[teo]{Lemma}}{}
\@ifundefined{prop}{\newtheorem{prop}[teo]{Proposition}}{}
\@ifundefined{remark}{\newtheorem{remark}[teo]{Remark}}{}
\makeatother
\DeclareMathOperator{\SO}{SO}
\DeclareMathOperator{\dist}{dist}
\renewcommand{\r}{\mathbb{R}}
\title[On the hierarchy of plate models for a singularly perturbed ]{On the hierarchy of plate models for a singularly perturbed multi-well nonlinear elastic energy}
\author{Edoardo Giovanni Tolotti}
\date{Source: arXiv:2501.11443v2 (CC BY 4.0)}
\begin{document}
\maketitle

\noindent\emph{Source and licence.} On the hierarchy of plate models for a singularly perturbed multi-well nonlinear elastic energy. Version arXiv:2501.11443v2 (CC BY 4.0), distributed under the Creative Commons Attribution 4.0 International licence, as recorded in the arXiv licence metadata for this exact version and shown on the abstract page https://arxiv.org/abs/2501.11443v2. Full rights notice: licenses/arxiv-2501.11443-CC-BY-4.0.txt.\par\smallskip

\noindent\textit{Editorial scope.} Counted unit: the Lemma whose source label is \texttt{lemma:compactness\_forces} in the article (source0.tex lines 1464--1470, source label \texttt{lemma:compactness\_forces}), delivered together with the article's own complete original proof of it (source0.tex lines 1472--1567, one \texttt{proof} environment of 96 lines). No mathematics is written, repaired or re-derived: statement and proof are the article's own text line for line. Required context retained uncounted: item:lower\_bound\_W (lines 214--219); eq:convergence\_forces (lines 419--421); prop:compactness (lines 695--712); prop:approximated\_rotation (lines 799--816); remark:approximated\_rotation\_Lr (lines 820--822). Every definition, display and label of those lines is retained unchanged. Omitted from this package and not redistributed: 1--213, 220--418, 422--694, 713--798, 817--819, 823--1463, 1471--1471, 1568--1915. Nothing inside a retained excerpt is omitted or altered: every retained body is delivered exactly as the article's lines are written, so no omission receipt of any kind accompanies this package. Citations: the retained excerpts carry no citation command at all, so no bibliography entry of the article is transcribed and no reference is redistributed. Reference adaptation: none was needed and none was made; every reference inside the delivered unit resolves inside it, so no unresolved reference or citation is produced. Printed slips kept literal: no printed word, symbol, hypothesis or number of the counted statement or of its proof was altered. Layout and class: the article's own class is \texttt{article} with packages including geometry, xcolor, hyperref, etoolbox, mathtools, mathrsfs, amsthm, cleveref, amsfonts, amssymb, autonum, biblatex, inputenc, enumitem; the wrapper is the lane's pinned \texttt{amsart} editorial wrapper at 10pt on A4, which restates the theorem-like environments the retained text needs and adds the article's own macro definitions that the retained text needs (\texttt{\textbackslash SO}, \texttt{\textbackslash dist}, \texttt{\textbackslash r}), with extra wrapper packages (bm, bbm, dsfont, xspace, cancel, mathtools, cleveref, enumitem). The delivered numbering is the unit's own. Assets: no retained span contains a figure environment, a graphics-inclusion command, a PDF-page inclusion, a table or a TikZ picture; no image file is copied into this package, the package declares no asset dependency and no image file is redistributed with it. Licence: version arXiv:2501.11443v2 of this article is distributed under the Creative Commons Attribution 4.0 International licence, as recorded in the arXiv licence metadata for this exact version and shown on the abstract page https://arxiv.org/abs/2501.11443v2. A full rights notice is delivered at licenses/arxiv-2501.11443-CC-BY-4.0.txt. Characters: every retained excerpt is pure ASCII, so the delivered engine, XeLaTeX with its default Latin Modern text font, reports no missing character.
\begin{enumerate}[start=1,label={(W\arabic*)}]
    \item $W(F) = 0 \iff F \in K$, \label{item:minimization_W}
    \item $W$ is $C^2$ in a neighborhood of $K$, \label{item:regularity_W}
    \item $W$ is frame indifferent, i.e., $W(RF) = W(F)$ for every $R \in \SO(3)$ and for every $F \in \r^{3 \times 3}$, \label{item:frame_indifference}
    \item $W(F) \geq C f_q(\dist(F, K))$ for every $F \in \r^{3 \times 3}$, \label{item:lower_bound_W}
\end{enumerate}

\begin{equation}\label{eq:convergence_forces}
	\dfrac{1}{h^{\gamma+1}}f_h \to f \quad \text{in} \, \, L^{q'}(S; \r^3) \,.
\end{equation}

\begin{prop}\label{prop:compactness}
    Let $\alpha \geq 2$. Let $(y_h) \subset W^{1,2}(\Omega; \r^3)$ be a sequence such that
    \begin{equation}
        \lim_{h \to 0} h^\alpha E^\alpha_h(y_h) = 0 \label{eq:hypothesis_compactness}\,.
    \end{equation}
    Then for $h \ll 1$ there are an index $i_h \in \{1, \dots, l\}$ and two constants $\delta, C(\delta) > 0$ such that
    \begin{align}
        & \int_{\Omega_{h}^{i_{h}}} \dist^2(\nabla_h y_h, K_{i_h}) \, dx \leq C(\delta)h^\alpha E^\alpha_h(y_h) \label{eq:compactness_near_well} \,,\\
        & \int_{\Omega \backslash \Omega_h^{i_h}} \dist^2(\nabla_h y_h, K_{i_h}) \, dx \leq C(\delta)h^\alpha[{(E^\alpha_h(y_h))}^{\theta} + E^\alpha_h(y_h)] \label{eq:compactness_far_well}\,,
    \end{align}
    where $\Omega_h^{i_h} = \{x \in \Omega \, : \, \dist(\nabla_h y_h(x), K_{i_h}) \leq \delta \}$ and
    \begin{equation}\label{eq:theta}
        \theta = \begin{cases}
            \frac{3}{2} \quad & \text{if} \, \, q = 2 \,,\\
            \frac{5}{3} \quad & \text{otherwise} \,.
        \end{cases}
    \end{equation}
\end{prop}

\begin{prop}\label{prop:approximated_rotation}
    Let $(y_h) \subset W^{1,2}(\Omega; \r^3)$ and let $j \in \{1, \dots, l\}$. Define
    \[
        D_{h, j} = \| \dist(\nabla_h y_h, K_j) \|_{L^2(\Omega)}\,.
    \]
    There are two maps $R_h \in L^\infty(S; \SO(3))$ and $\tilde R_h \in W^{1,2}(S; \r^{3 \times 3}) \cap L^\infty(S; \r^{3 \times 3})$ such that
    \begin{enumerate}[start=1,label={(R\arabic*)}]
        \item $\|\nabla_h y_h - R_h U_j\|_{L^2(\Omega)} \leq C D_{h, j}$\label{item:approximated_rotation_convergence_scaled_gradient},
        \item $\|\nabla \tilde R_h\|_{L^2(S)} \leq C h^{-1}D_{h, j}$\label{item:approximated_rotation_convergence_gradient},
        \item $\|\tilde R_h - R_h\|_{L^2(S)} \leq C D_{h, j}$,
        \item $\|\tilde R_h - R_h\|_{L^\infty(S)} \leq C h^{-1}D_{h, j}$.
    \end{enumerate}
    Moreover, there exists a constant rotation $Q_h \in \SO(3)$ such that
    \begin{equation}
        \|R_h - Q_h\|_{L^2(S)} \leq Ch^{-1}D_{h, j} \,.\label{eq:approximated_constant_rotation}
    \end{equation}
    Finally, if $h^{-1}D_{h, j} \to 0$, then for $h \ll 1$ we can choose $\tilde R_h = R_h$.
\end{prop}

\begin{remark}\label{remark:approximated_rotation_Lr}
	If $ r > 1 $, all the results of \Cref{prop:approximated_rotation} hold with the $L^2$ norm replaced by the $L^r$ norm and the factor $h^{-1}$ replaced by $h^{-\frac{2}{r}}$.
\end{remark}

\begin{lemma}\label{lemma:compactness_forces}
	Let $\alpha \geq 2$ and $q > 1$. Suppose that $(y_h) \subset W^{1,2}(\Omega; \r^3)$ is a sequence of deformations that are quasi-minimizers for $J_h^\alpha$, that is,
	\[
	\lim_{h \to 0} \, (J_h^\alpha(y_h) - \inf J_h^\alpha) =0 \,.
	\]
	Then, $E_h^\alpha(y_h) \leq C$ for every $h > 0$.
\end{lemma}

\begin{proof}
	Firstly, we will prove that $h^\alpha E_h^\alpha(y_h) \to 0$. Fix $j \in \{1, \dots, l\}$ and let $Q_h$, $R_h$ be, respectively, the rotation and the map given by in \Cref{prop:approximated_rotation}. Define
	\[
	\tilde y_h = y_h  - Q_h U_j\begin{pmatrix}
	x' \\ h x_3
	\end{pmatrix} - c_h \,,
	\] 
	where
	\[
	c_h = \dfrac{1}{|\Omega|}\int_\Omega \left(y_h - Q_h U_j \begin{pmatrix}
	x' \\ hx_3
	\end{pmatrix}\right) \, dx \,.
	\]
	Using the test deformation
	\[
	(x', x_3) \mapsto U_j \begin{pmatrix}
	x' \\ hx_3
	\end{pmatrix}
	\]
	and the assumption \eqref{eq:convergence_forces}, we get $\inf_y J_h^\alpha(y) \leq Ch^{1-\gamma}$, so that
	\[
		J_h^\alpha(y_h) - C h^{1-\gamma} \leq J_h^\alpha(y_h) - \inf_y J_h^\alpha(y) \leq C \implies J_h^\alpha(y_h) \leq Ch^{1-\gamma} \,.
	\]
	We have
	\begin{equation}\label{eq:dim_forze_eq_1}
	\begin{aligned}
	h^\alpha E_h^\alpha(y_h) & = h^\alpha J_h^\alpha(y_h) + \int_\Omega f_h \cdot \tilde y_h \, dx  + \int_S f_h \cdot Q_hU_j \begin{pmatrix}
	x' \\ 0
	\end{pmatrix} \, dx \\
	& \leq Ch^{\gamma+1} + Ch^{\gamma+1}\|\nabla_h \tilde y_h\|_{L^q} \,.
	\end{aligned}
	\end{equation}
	Consider the set
	\[
	\tilde \Omega_h = \{x \in \Omega \colon \dist(\nabla_h y_h, K) \geq 1\}\,.
	\]
	By \Cref{prop:approximated_rotation} and \Cref{remark:approximated_rotation_Lr}, we deduce
	\begin{equation}\label{eq:dim_forze_eq_2}
		\begin{aligned}
		\|\nabla_h \tilde y_h\|_{L^q}^q & = \int_\Omega |\nabla_h y_h - Q_hU_j|^q \, dx \leq Ch^{-2} \int_\Omega \dist^q(\nabla_h y_h, K_j) \, dx\\
		& \leq Ch^{-2}\int_\Omega \dist^q(\nabla_h y_h, K) \, dx + Ch^{-2}\\
		& \leq Ch^{-2} \int_{\tilde \Omega_h} \dist^q(\nabla_h y_h, K) \, dx + Ch^{-2}\\
		& \leq Ch^{\alpha-2} E_h^\alpha(y_h) + Ch^{-2} \,,
		\end{aligned}
	\end{equation}
	where we have used \ref{item:lower_bound_W}. 
	Combining \eqref{eq:dim_forze_eq_1}-\eqref{eq:dim_forze_eq_2} we get
	\[
	h^\alpha E_h^\alpha(y_h) \leq Ch^{\gamma + 1} + Ch^{\gamma+1 - \frac{2}{q}}(h^\alpha E_h^\alpha(y_h))^{\frac{1}{q}}  +Ch^{\gamma + 1 - \frac{2}{q}} \,.
	\]
	Recalling that $q > 1$, by Young's inequality we deduce that
	\[
	h^\alpha E_h^\alpha(y_h) \leq Ch^{\gamma +1} + Ch^{q'(\gamma + 1 - \frac{2}{q})} + Ch^{\gamma + 1 - \frac{2}{q}} \to 0\,.
	\]
	Hence, we can apply \Cref{prop:compactness} and deduce that, at least along a subsequence, there is an index $j_0 \in \{1, \dots, l\}$ such that for $h \ll 1$
	\[
	\int_\Omega \dist^{2}(\nabla_h y_h, K_{j_0}) \,, dx \leq h^{\alpha}[E_h^{\alpha}(y_h) + (E_h^{\alpha}(y_h))^\theta] \,.
	\]

	We now show that $E_h^\alpha(y_h) \leq C$. To simplify the exposition, we suppose that $j_0 = j$. By \Cref{prop:approximated_rotation} we have
	\begin{equation}\label{eq:bound_yh_energy}
		\begin{aligned}
			\|\tilde y_h\|_{L^2}^2 & \leq C \|\nabla \tilde y_h\|_{L^2}^2 \leq Ch^{-2} \|\dist(\nabla_h y_h, K_j)\|^2_{L^2}\\
			& \leq Ch^{\alpha-2}[E_h^{\alpha}(y_h) + (E_h^{\alpha}(y_h))^\theta] \,.
		\end{aligned}
	\end{equation}
	Pick $\bar R_h U_{k_h} \in \mathcal{M}_h$ and define
	\[
	\tilde J^\alpha_h(y) := E^{\alpha}_h(y) - \dfrac{1}{h^\alpha} \int_\Omega f_h \cdot y \, dx + \dfrac{1}{h^\alpha} \int_S f_h \cdot \bar R_h U_{k_h} \begin{pmatrix}
	x' \\ 0
	\end{pmatrix} \, dx \,.
	\]
	Note that $y_h$ are quasi-minimizers for the functional $\tilde J_h^\alpha$. Moreover, using the test deformation
	\[
		(x', x_3) \mapsto \bar R_h U_{k_h}\begin{pmatrix}
		x' \\ hx_3
		\end{pmatrix} \,,
	\]
	we get $\inf_y \tilde J_h^\alpha(y) \leq 0$, so that
	\[
		\tilde J_h^\alpha(y_h) \leq \tilde J_h^{\alpha}(y_h) - \inf_y \tilde J_h^\alpha(y) \leq C \,.
	\]
	Hence,
	\begin{align}
	E_h^\alpha(y_h) & = \tilde J^\alpha_h(y_h) + \dfrac{1}{h^\alpha} \int_\Omega f_h \cdot y_h \, dx - \dfrac{1}{h^\alpha} \int_S f_h \cdot \bar R_h U_{k_h} \begin{pmatrix}
	x' \\ 0
	\end{pmatrix} \, dx' \\
	& = \tilde J^\alpha_h(y_h) + \dfrac{1}{h^\alpha} \int_\Omega f_h \cdot \tilde y_h \, dx + \dfrac{1}{h^\alpha} \int_S f_h \cdot (Q_h U_j  - \bar R_h U_{k_h}) \begin{pmatrix}
	x' \\ 0
	\end{pmatrix} \, dx'\\
	& \leq C + Ch^{-\gamma -1}\|f_h\|_{L^2}(E_h^{\alpha}(y_h) + (E_h^\alpha(y_h))^\theta )^{\frac{1}{2}}\\
	& \leq C + C(E_h^\alpha(y_h))^\frac{1}{2} + C(E_h^\alpha(y_h))^\frac{\theta}{2}\,,
	\end{align}
	where we have used \eqref{eq:bound_yh_energy} and the definition of $\mathcal{M}_h$.
	Since $0 < \frac{\theta}{2} < 1$, by an application of Young's inequality we conclude.
\end{proof}

\end{document}
